\documentclass[../../../include/open-logic-section]{subfiles}

\begin{document}

\olfileid{sth}{choice}{hartogs}
\olsection{القابلية للمقارنة ولمّة هارتوغس}

هذا وجه الأمر النافع؛ وأما وجهه الآخر، فإن الحال من غير الاختيار يختل
اختلالًا شديدًا. وتظهر العلة في النتيجة الحسنة الآتية
لـ~\cite{Hartogs1915}.

\begin{lem}[\emph{في $\ZF$}]\ollabel{HartogsLemma}
لكل مجموعة~$A$، يوجد عدد ترتيبي~$\alpha$ بحيث
$\cardnless{\alpha}{A}$.
\end{lem}

\begin{proof}
إذا كان~$B \subseteq A$ و$R \subseteq B^2$، فإن
$\tuple{B, R} \subseteq V_{\setrank{A}+4}$، وذلك بمقتضى
\olref[ord-arithmetic][using-addition]{rankcomputation}. فانظر،
باستعمال الفصل، إلى المجموعة:
\[
	C = \Setabs{\tuple{B, R} \in V_{\setrank{A}+5}}{B\subseteq A
	\text{ و$\tuple{B, R}$ ترتيب جيد}}
\]
ثم كوّن، بالاستبدال وبالاستعانة بـ
\olref[ordinals][ordtype]{thmOrdinalRepresentation}، المجموعة:
\[
	\alpha = \Setabs{\ordtype{B, R}}{\tuple{B, R} \in C}.
\]
وهذه~$\alpha$ عدد ترتيبي بموجب
\olref[ordinals][basic]{corordtransitiveord}، لأنها مجموعة متعدية من
الأعداد الترتيبية. ووجه ذلك أنه إذا كان~$\gamma \in \beta \in \alpha$،
فإن~$\beta = \ordtype{B, R}$ لبعض~$B \subseteq A$، وحينئذ
$\gamma = \ordtype{B_b, R_b}$ لبعض~$b \in B$، على ما في
\olref[ordinals][iso]{wellordinitialsegment}؛ ومن ثم
$\gamma \in \alpha$.

ونطلب التناقض، فنفرض وجود !!a{injection}
$f \colon \alpha \to A$. وعندئذ ضع:
\begin{align*}
	B &= \ran{f}\\
	R &= \Setabs{\tuple{f(\gamma), f(\delta)} \in A \times A}
	  {\gamma,\delta \in \alpha \land \gamma \in \delta}.
\end{align*}
فواضح أن~$\alpha = \ordtype{B, R}$ وأن
$\tuple{B, R} \in C$. فيلزم~$\alpha \in \alpha$، وهذا محال.
% تصحيح مصدر (0595-I1): فُصل متغيرا الربط gamma وdelta عن العدد الترتيبي alpha وقُيدا بمجال f.
\end{proof}

وينشأ عن ذلك حكم عميق.

\begin{thm}[\emph{في $\ZF$}]
الادعاءان الآتيان متكافئان:
\begin{enumerate}
	\item\ollabel{equivwo} مسلّمة الترتيب الجيد؛
	\item\ollabel{equivcompare} إما~$\cardle{A}{B}$ أو
	$\cardle{B}{A}$، لأي مجموعتين~$A$ و$B$.
\end{enumerate}
\end{thm}

\begin{proof}
\emph{\olref{equivwo} $\Rightarrow$ \olref{equivcompare}.} خذ~$A$
و$B$. من~\olref{equivwo} يوجد ترتيبان جيدان
$\tuple{A, R}$ و$\tuple{B, S}$. ومن
\olref[ordinals][ordtype]{thmOrdinalRepresentation} يوجد تماثلان
$f \colon \alpha \to \tuple{A, R}$ و
$g \colon \beta \to \tuple{B, S}$. وبحسب
\olref[sth][ordinals][basic]{ordinalsaresubsets}، إما
$\alpha \subseteq \beta$ أو~$\beta \subseteq \alpha$. فإن كان
$\alpha \subseteq \beta$، كانت
$\comp{f^{-1}}{g} \colon A \to B$ !!a{injection}، فينتج
$\cardle{A}{B}$؛ وإن كان~$\beta \subseteq \alpha$، كان بالمثل
$\cardle{B}{A}$.

\emph{\olref{equivcompare} $\Rightarrow$ \olref{equivwo}.} خذ~$A$؛
فبموجب~\olref{HartogsLemma} يوجد عدد ترتيبي~$\beta$ بحيث
$\cardnless{\beta}{A}$. ثم يعطي~\olref{equivcompare} أن
$\cardle{A}{\beta}$. فتوجد !!{injection}
$f \colon A \to \beta$، وبهذا الحقن نرتب عناصر~$A$ ترتيبًا جيدًا،
بأن نعرف الترتيب
$\Setabs{\tuple{a, b} \in A \times A}{f(a) \in f(b)}$.
\end{proof}
\noindent
ومن هذا تتبع نتيجة قريبة: إذا بطل الترتيب الجيد، كانت بعض المجموعات غير
قابلة للمقارنة أصلًا من جهة الحجم، واضطرب الحساب الكاردينالي فيما وراء
المتناهي. فلا يبقى لنا مثلًا أن نقول، إذا كانت~$A$ و$B$ غير منتهيتين،
إن~$\cardeq{A \disjointsum B}{M}$ و$\cardeq{A \times B}{M}$، حيث~$M$
كبرى~$A$ و$B$؛ وانظر
\olref[card-arithmetic][simp]{cardplustimesmax}. والعلة ظاهرة: متى لم
يمكننا مقارنة حجم~$A$ بحجم~$B$، لم يستقم السؤال عن الأكبر منهما.

\end{document}
